\documentclass[a4paper,12pt]{article}

% Paquetes
\usepackage[spanish]{babel}
\usepackage[utf8]{inputenc}
\usepackage[T1]{fontenc}
\usepackage{amsmath, amssymb}
\usepackage{graphicx}
\usepackage{siunitx}
\usepackage{geometry}
\usepackage{float}

\geometry{margin=2.5cm}

\title{Práctica I: Instrumentación electrónica básica}
\author{Nombre Apellido}
\date{\today}

\begin{document}

\maketitle

\tableofcontents
\newpage


\section{Intrumentación electrónica}

%----------------------------------------
\subsection{Objetivo}

%----------------------------------------
\subsection{Introducción}

%----------------------------------------
\subsection{Resistencias individuales}


\begin{table}[H]
\centering
\begin{tabular}{cccc}
\hline
Resistencia & Fabricación (k$\Omega$) & Polímetro (k$\Omega$) & Discrepancia (\%) \\
\hline
1 & 0.56 & 0.558 & 0.36 \\
2 & 10 & 10.02 & 0.20 \\
3 & 0.27 & 0.267 & 1.11 \\
4 & 1.6 & 1.606 & 0.38 \\
\hline
\end{tabular}
\caption{Comparación entre valores nominales y medidos de resistencias}
\end{table}


%----------------------------------------
\subsection{Ley de Ohm}

%----------------------------------------
\subsection{Resistencia en un circuito}

%----------------------------------------
\subsection{Leyes de Kirchhoff}

%----------------------------------------
\subsection{Teoremas de Thevenin y Norton}

%----------------------------------------
\subsection{Teorema de transferencia de máxima potencia}

%----------------------------------------
\subsection{Conclusiones}

%----------------------------------------
\subsection{Cuestiones}

%----------------------------------------
\begin{thebibliography}{9}

\bibitem{ref1}
Referencia aquí

\end{thebibliography}

\end{document}